% ============================================================
% APÉNDICE (ubicar al final del documento)
% ============================================================
\appendix
\refstepcounter{section}
\section*{\centering Apéndice \thesection}\label{apx:cuestionario}
\addcontentsline{toc}{section}{Apéndice \thesection}

\subsubsection*{Instrumento: Cuestionario WebSIS}

Estimado/a estudiante: Este cuestionario es anónimo y tiene como finalidad mejorar la plataforma WebSIS. Por favor, responda con sinceridad. Duración estimada: 5 minutos.

\paragraph{Sección A: Datos del encuestado}
\begin{itemize}
  \item \textbf{Sexo:} [ ] Femenino \quad [ ] Masculino
  \item \textbf{Edad:} [ ] 17--19 años \quad [ ] 20--22 años \quad [ ] 23--25 años \quad [ ] Más de 25 años
  \item \textbf{Estado civil:} [ ] Soltero/a \quad [ ] Casado/a \quad [ ] Unión libre \quad [ ] Divorciado/a
  \item \textbf{¿Actualmente trabajas?} [ ] Sí \quad [ ] No
  \item \textbf{Facultad a la que pertenece:} \rule{6cm}{0.4pt}
  \item \textbf{Semestre que cursa actualmente:} [ ] 1--2 \quad [ ] 3--4 \quad [ ] 5--6 \quad [ ] 7 o más
\end{itemize}

\paragraph{Sección B: Uso de WebSIS}
\begin{itemize}
  \item \textbf{B1. ¿Con qué frecuencia utilizas WebSIS?} [ ] Diariamente \quad [ ] Varias veces por semana \quad [ ] Solo en períodos de inscripción/notas \quad [ ] Raramente
  \item \textbf{B2. ¿Para qué tareas utilizas WebSIS principalmente?} Puedes marcar más de una: [ ] Inscripción a materias \quad [ ] Consulta de notas \quad [ ] Consulta de horarios \quad [ ] Situación académica \quad [ ] Otra: \rule{2cm}{0.4pt}
  \item \textbf{B3. ¿Desde qué dispositivo accedes habitualmente a WebSIS?} [ ] Computadora de escritorio \quad [ ] Laptop \quad [ ] Teléfono celular \quad [ ] Tablet
  \item \textbf{B4. ¿Desde qué lugar accedes normalmente a WebSIS?} \par [ ] Casa \quad [ ] Universidad/laboratorio/Wi-Fi \quad [ ] Cybercafé \quad [ ] Trabajo
\end{itemize}

\paragraph{Sección C: Frustraciones y barreras}
\begin{itemize}
  \item \textbf{C1. ¿Qué tan difícil te resulta usar WebSIS en general?} [ ] Muy fácil \quad [ ] Fácil \quad [ ] Neutral \quad [ ] Difícil \quad [ ] Muy difícil
  \item \textbf{C2. ¿Qué tan difícil te resulta la inscripción a materias?} [ ] Muy fácil \quad [ ] Fácil \quad [ ] Neutral \quad [ ] Difícil \quad [ ] Muy difícil
  \item \textbf{C3. ¿Qué tan difícil te resulta consultar tus notas?} [ ] Muy fácil \quad [ ] Fácil \quad [ ] Neutral \quad [ ] Difícil \quad [ ] Muy difícil
  \item \textbf{C4. ¿Con qué frecuencia necesitas ayuda de otra persona para completar una tarea en WebSIS?} [ ] Nunca \quad [ ] A veces \quad [ ] Frecuentemente \quad [ ] Siempre
  \item \textbf{C5. ¿Alguna vez no pudiste completar una tarea importante, por ejemplo, inscripción a tiempo?} [ ] Sí \quad [ ] No. Si tu respuesta es Sí, ¿qué ocurrió?: \rule{4cm}{0.4pt}
  \item \textbf{C6. ¿Cuál es el problema más frecuente que encuentras?} [ ] No sé dónde encontrar lo que busco \quad [ ] La interfaz es confusa \quad [ ] El sistema es lento o falla \quad [ ] No entiendo los mensajes de error \quad [ ] Los pasos a seguir no son obvios \quad [ ] Otro: \rule{2cm}{0.4pt}
  \item \textbf{C7. En una escala del 1 al 5, ¿qué tan satisfecho/a estás con WebSIS?} 1[ ] \quad 2[ ] \quad 3[ ] \quad 4[ ] \quad 5[ ] \quad (1 = Muy insatisfecho, 5 = Muy satisfecho)
  \item \textbf{C8. ¿Qué mejora considerarías más importante?} [ ] Mejor organización del menú \quad [ ] Mensajes de error más claros \quad [ ] Guía o tutorial integrado \quad [ ] Diseño más moderno \quad [ ] Versión móvil optimizada \quad [ ] Otro: \rule{2cm}{0.4pt}
  \item \textbf{C9. Cuando buscas una función en WebSIS, ¿dónde esperarías encontrar primero opciones como inscripción, notas o horario?} [ ] En un menú principal por tareas \quad [ ] En una sección por trámites \quad [ ] En un panel con accesos directos \quad [ ] No lo tengo claro
  \item \textbf{C10. ¿Qué nombre te resultaría más natural para agrupar funciones como notas, Kardex y situación académica?} [ ] Rendimiento académico \quad [ ] Historial académico \quad [ ] Mi avance académico \quad [ ] No sabría cómo llamarlo
  \item \textbf{C11. Al iniciar una tarea importante, como la inscripción, ¿qué esperarías ver primero?} [ ] Los pasos a seguir \quad [ ] Mi estado habilitado/no habilitado \quad [ ] Las materias disponibles \quad [ ] Un resumen general del proceso
  \item \textbf{C12. Si el sistema muestra un error o bloqueo, ¿qué esperarías que haga para ayudarte?} [ ] Explicar qué ocurrió \quad [ ] Indicar cómo solucionarlo \quad [ ] Señalar en qué paso estoy \quad [ ] Permitir corregir sin empezar de nuevo
\end{itemize}

\paragraph{Sección D: Experiencia emocional}

\begin{itemize}
  \item \textbf{D1. Cuando necesitas usar WebSIS en un período de inscripción o publicación de notas, ¿cómo te sientes antes de ingresar al sistema?} [ ] Tranquilo/a \quad [ ] Algo ansioso/a \quad [ ] Muy ansioso/a \quad [ ] Indiferente
  \item \textbf{D2. ¿Cómo te sientes generalmente después de completar una tarea en WebSIS?} [ ] Satisfecho/a y seguro/a \quad [ ] Aliviado/a pero con dudas sobre si lo hice bien \quad [ ] Frustrado/a \quad [ ] Indiferente
  \item \textbf{D3. Si no puedes completar una tarea, ¿cuál es tu reacción más frecuente?} [ ] Busco ayuda de un compañero \quad [ ] Lo intento de nuevo más tarde \quad [ ] Me rindo
  \item \textbf{D4. ¿Sientes que el sistema confía en ti como usuario, es decir, te da información suficiente para actuar sin ayuda?} [ ] Sí, siempre \quad [ ] A veces \quad [ ] Rara vez \quad [ ] No, casi nunca
  \item \textbf{D5. ¿Qué palabra describe mejor tu relación con WebSIS?} Elige una: [ ] Frustración \quad [ ] Tolerancia \quad [ ] Confianza \quad [ ] Indiferencia \quad [ ] Dependencia
\end{itemize}
